\documentclass[pre,twocolumn,superscriptaddress,nofootinbib]{revtex4-2}
\usepackage{graphicx}
\usepackage{amsmath}
\usepackage{booktabs}
\usepackage[hidelinks]{hyperref}

\begin{document}

\title{Why a square-counting graph energy shows no pull at a distance, and what a pull would need}

\author{Emily Smith}
\email{emilysmithcreate@gmail.com}
\affiliation{Independent researcher, Charlottesville, Virginia, USA}

\date{DRAFT, 25 September 2026}

\begin{abstract}
The author's hypothesis holds that gravity is the wish of curved space to refold: that small curled leftovers of the
change that made space pull on each other through the arrangement around them. We ask whether the family of graph
energies obtained from combinatorial quantum gravity (CQG) by raising the coefficient $\lambda$ of its local term above
1 can produce such a pull. At fixed wiring, exact results say it does not. The energy is a sum over links, so two
defects that share no square are exactly additive; counting symmetries, in the cases computed, prefers symmetric
placements, not near ones; and flat space is gapped in the extreme: every single move out of it costs 32, 64 or 128
units with four, six or eight links per point, at every $\lambda$. That a medium with such a gap carries a force only
over a short range is argued from the published literature, by analogy, not derived for this model; a force with the
wiring free to rearrange has not been measured, and with six links the relic such a force would act on has not been
found. We then state what the published work says a pull needs (a mode with no gap, each relic a source for it in
proportion to its energy, and a mode of the kind that makes like sources attract), and compute exactly what the
simplest such mode a graph carries, a field on its points, would add: its fluctuations give a Casimir-type attraction
under the field reading, a free-energy term falling roughly as $r^{-4}$ in two directions and $r^{-6}$ in three, worth
about 0.02 of the coupling at contact in two directions and 0.002 in three; with relics as its sources it gives a
potential falling as $1/r$ in three directions, Newton's shape, carried by the emergent wiring; and the field barely
changes which arrangement is favored. That rule would put the pull in, not derive it; whether to adopt it is the next
decision.
\end{abstract}

\maketitle

\section{Introduction}

The companion papers \cite{paper1,paper2} study a curled torus that opens into flat space, the energy it releases and the
small curled leftover it leaves behind, the relic. The author's framework reads gravity as the refolding drive of curved
space: relics, being curled, should draw one another together through the arrangement between them. That makes a
definite question of the model: do two relics in flat space pull on each other, and if so, how does the pull fall with
distance? We answer it for the model as it stands, say why the answer is what it is, and say what would change it.

The energy is that of CQG \cite{T25} with the coefficient of its local term raised: $H = 16(D(D-1)N/2 - S) + 4\lambda X$
on $2D$-regular bipartite graphs, $S$ the number of 4-cycles (squares) and $X$ the surplus squares on over-full links.
Only $\lambda = 1$ is CQG; everything run here is at $\lambda > 1$, next to it. The exact statements hold at every
$\lambda$, including 1. The six- and eight-link energies are derived by us from the published definitions; the
six-link kernel has not reproduced the published six-link curve, and there is no published eight-link curve, so every
statement here at six and eight links is about our family and its kernel.

\section{Three exact facts}

\emph{Additivity.} $H$ is a sum of terms on links, each depending on the squares through that link. Two defects that
share no square therefore cost exactly the sum of their separate costs, at every separation; an attraction appears only
at contact (16 with four links, 64 with six), where they share squares. We checked it by construction for local defects at
every separation on flat tori with four and with six links.

\emph{No counting pull at fixed wiring.} With the points treated as interchangeable, as the author's framework treats
them, each placement of two identical defects is weighted by its number of symmetries. With four links, for one simple
defect on one torus, that number is one or two at ordinary separations and four only at exactly opposite placements;
with six links it is 8 at generic separations, 16 exactly opposite, 32 along the defect's own axis at any separation
and 4 on a diagonal. That is a preference for symmetric placements, not one that grows as the defects approach: not a
pull.

\emph{A gapped vacuum.} Flat space is the exact ground state for $\lambda \ge 1$, and the cheapest single move out of it
costs 32, 64 and 128 with four, six and eight links, the same at every $\lambda$ (listed over every switch with four
links; with six and eight, by a search over nearby partners that was checked against the full search on one six-link
torus). At the couplings where curled arrangements are stuck for now ($g \lesssim 2$), a cold flat region accepts
essentially no move at all; the relic itself anneals away within 500 to 80,500 sweeps at $g = 1.0$ to 1.5 and lasts
only where its surroundings stay very cold \cite{paper2}. The relic is a dip with four links (it costs
$24\lambda - 16$ and every single move out of it costs energy); with six links, a search at $\lambda = 1.02$ on the flat
$6 \times 6 \times 6$ torus, over the 9,049 curled arrangements made by one or two switches within three steps of one
point, testing the 30 lowest in energy, found none that is a dip, so the object a pull would act on is not yet known to
exist there.

\section{Why the facts are one fact}

A force between two objects in a medium is carried by the medium's fluctuations, and its range is the range of their
correlations \cite{KG99}. A gapped system has correlations that decay exponentially \cite{HK06}; a massive mediator gives
Yukawa's $e^{-mr}/r$ \cite{Tong}. The model's author writes that both phases of the model are governed by a finite
correlation length, which diverges only at the critical point \cite{T25}. Here ``gapped'' means that every single move
out of flat space costs at least a fixed energy; the theorem of Ref.~\cite{HK06} is for quantum lattice Hamiltonians,
and carrying it over to this classical ensemble, defined by its moves, is an analogy we have not proved. A vacuum whose
every excitation costs 32 or more is the extreme case: at fixed wiring there is nothing to carry a force (the first two
facts), and with the wiring free there is almost nothing that moves. So, by this argument, the absence of a pull
between point defects is what this energy should give in its geometric phase, not a failure to look hard enough; a
force with the wiring free to rearrange, which would be entropic, has not been measured. One long-range force does
exist in this family, of another kind: because two orders differ by a fixed energy per point, a boundary between them
is pushed with a constant force, the force that drives the front of Ref.~\cite{paper1}. It acts between regions of
different order, not between separated lumps, and it is not gravity.

We had also measured, in the published model ($\lambda = 1$), that the curvature correlation falls to zero within two
steps at every coupling and every size up to 676 points; the same measurement did not reproduce the model's
published correlation length.
That is consistent with a gap but does not show one: in Einstein's theory, which has a massless graviton, the curvature
correlator is a contact term at tree level, and it is the volume correlator that carries the massless mode \cite{LR25}.
The gap is shown by the walls.

\section{What a pull needs}

The published work names three things, all needed together. (i) A mode with no gap, so that the medium's correlations
fall as a power, not an exponential. The robust classical route is a hard local constraint with many ground states, as
in dimer models on bipartite lattices, where counting alone gives defects an inverse-square force \cite{HKMS03,Henley};
the quantum route is a gauge structure \cite{GW09}. (ii) Each relic a source for that mode in proportion to its energy,
a charge whose flux through every surrounding surface is the same; without it a gapless medium gives only the fast
Casimir laws that stiff inclusions feel in a membrane \cite{GBP93,KG99}. (iii) A mode of the kind that makes like sources
attract: a scalar, as in Nordstr\"om's theory \cite{Giulini}, or the tensor of Einstein's; the vector modes that hard
constraints produce make like defects repel \cite{Henley}, as electric charges do. In three large directions the three
together give Newton's $1/r$ potential. The one discrete random-geometry model we have found reported to show a
Newtonian pull, four-dimensional
dynamical triangulations, did so with test particles moving on a geometry tuned near a transition \cite{dBS97,DLSU21}.

\section{The simplest candidate, computed}

A graph carries one mode with no gap for free: a field $\phi$ on its points with energy $\frac{\kappa}{2}\sum_{\rm
links}(\phi_u - \phi_v)^2$, whose propagator is the graph's own Green's function and so follows the emergent geometry.
Integrated out exactly, it adds to the free energy of a wiring $\frac{g}{2}\ln\det{}'L$ (with $L$ the graph Laplacian;
$\det{}'L$ is $N$ times the number of spanning trees), and, if relics are sources $q$ for it, a cross term
$-q_1^{\mathsf T}L^{+}q_2$ between two relics. We evaluated both on fixed wirings holding two identical defects at every
even separation (Table~\ref{tab:pull}), each with its expectation written before the calculation. The defect is the
one-switch curled line of four, which is not a relic (with six links it is not a dip, and with four it is not the
relic of Ref.~\cite{paper2}, which adds a square where a single switch from flat space cannot), with a unit charge on
each of its points; for two identical defects that differs from a charge in proportion to energy only by a common factor. The
fluctuation term is lower when the defects are close, so under the field reading it pulls them together (if the
spanning trees were counted instead as hidden structure, the sign would reverse), but weakly and briefly, falling
roughly as $r^{-4}$ in two directions and $r^{-6}$ in three (local exponents 2.8 to 4.9 and 4.2 to 6.2 across the
separations of Table~\ref{tab:pull}) and worth about $0.024g$ in two directions and $0.0016g$ in three at contact: it
passes the test of sign and fails that of shape. The source term, a potential, grows as the relics approach, as the
lattice Green's function does, like $1/r$ in three directions, less the images of the periodic torus, and like
$-\ln r$ in two. Across the rungs of the ladder the field's own term favors a curled torus over flat space by
0.0144$g$ per point with four links, and by 0.0026$g$ to 0.0039$g$ per point per curled direction with six (one and
two curled): the same as lowering $\lambda$ by about $0.004g$ with four links and $0.001g$ or less with six, too little
to move the window in which the change happens at $\lambda = 1.25$.

\begin{table}[t]
\caption{Two identical defects (the one-switch curled line of four, not a relic) at separation $r$ along an axis, on
fixed wirings: the change in
$\ln$(spanning trees) and in the source cross term $q_1^{\mathsf T}L^{+}q_2$ (unit charge on each defect point), each
relative to the farthest placement. Exact.}
\label{tab:pull}
\begin{ruledtabular}
\begin{tabular}{lccc}
torus & $r$ & $\Delta\ln\tau$ & $\Delta\, q_1^{\mathsf T}L^{+}q_2$ \\
\hline
$24 \times 24$ & 2 & $-0.0478$ & 3.54 \\
 & 4 & $-0.0068$ & 1.83 \\
 & 6 & $-0.0014$ & 0.91 \\
 & 8 & $-0.00034$ & 0.38 \\
$16 \times 16 \times 16$ & 2 & $-0.0032$ & 0.369 \\
 & 4 & $-0.00017$ & 0.105 \\
 & 6 & $-0.000014$ & 0.022 \\
\end{tabular}
\end{ruledtabular}
\end{table}

\section{Discussion}

What this paper shows is narrow. In this family of energies at $\lambda > 1$, at fixed wiring, point defects that
share no square do not pull on each other through the energy, and counting symmetries gives no attraction that grows
as they approach, in the cases computed; the gap of flat space argues, by analogy with the published results, that
nothing longer-ranged is carried either, whatever the size of the runs. At $\lambda = 1$, which is CQG, we claim only
the exact facts. A force with the wiring free to rearrange has not been measured, and with six links the relic it
would act on has not been found. It does not show that the author's account of gravity is wrong; it argues that this
energy does not contain the ingredient that the published accounts of a long-range pull need. The simplest ingredient that would, a massless field on the points with relics as its sources,
gives Newton's shape, but it puts the pull in by hand; what would be new is that the field travels through the emergent
wiring rather than a fixed lattice, and can act back on it. Three tests follow if it is adopted, each to be
pre-registered with the author's prediction: whether two relics drift together in a sealed run; the shape of the pull in
three directions and its scaling with the relics' energies; and whether many relics clump under it and a dense clump
re-curls space, which is the author's mechanism for black holes \cite{paper4}. Whether the field could itself arise from
the loops of the graph, rather than be added, is the harder question behind all three.

\begin{acknowledgments}
Code, analysis and text were drafted with Claude (Anthropic) in conversation with the author, who directed and checked
the work; the literature summarized in Secs.~III and IV was read by AI agents from fetched copies and is to be checked
against the papers; no physicist has reviewed this. Every configuration, seed, result and pre-registration is in the
repository: \url{https://github.com/EmilySmithCreate/RecreationalPhysics} (ASSUMPTIONS O22, O23, O31, O32, O35, O50,
O51, O56, O58, O60, O61;
\texttt{scripts/exact\_tree\_count\_pull.py}, \texttt{scripts/exact\_tree\_count\_ladder.py}).
\end{acknowledgments}

\begin{thebibliography}{99}
\bibitem{paper1} E. Smith, ``The curled torus burps: front-driven decompactification between ordered states in a graph
model of emergent geometry'' (submitted to arXiv, 2026).
\bibitem{paper2} E. Smith, ``What the burp leaves behind: a relic defect of fixed size in a graph model of emergent
geometry, its number, its position, and how long it lasts,'' draft (2026).
\bibitem{paper4} E. Smith, ``Does concentrated energy re-curl space? Three sealed tests in a graph model of emergent
geometry, and a melt every time,'' draft (2026).
\bibitem{T25} C. A. Trugenberger, ``Networks as the fundamental constituents of the universe,'' arXiv:2512.17676.
\bibitem{KG99} M. Kardar and R. Golestanian, ``The `friction' of vacuum, and other fluctuation-induced forces,'' Rev. Mod.
Phys. \textbf{71}, 1233 (1999), arXiv:cond-mat/9711071.
\bibitem{HK06} M. B. Hastings and T. Koma, ``Spectral gap and exponential decay of correlations,'' arXiv:math-ph/0507008.
\bibitem{Tong} D. Tong, \emph{Quantum Field Theory}, lecture notes, University of Cambridge, Ch. 3.
\bibitem{LR25} J. Laiho and K. Ratliff, ``Euclidean correlation functions in quantum gravity,'' arXiv:2510.11888.
\bibitem{HKMS03} D. A. Huse, W. Krauth, R. Moessner and S. L. Sondhi, ``Coulomb and liquid dimer models in three
dimensions,'' arXiv:cond-mat/0305318.
\bibitem{Henley} C. L. Henley, ``The `Coulomb phase' in frustrated systems,'' arXiv:0912.4531.
\bibitem{GW09} Z.-C. Gu and X.-G. Wen, ``Emergence of helicity $\pm 2$ modes (gravitons) from qubit models,''
arXiv:0907.1203.
\bibitem{GBP93} M. Goulian, R. Bruinsma and P. Pincus, ``Long-range forces in heterogeneous fluid membranes,'' Europhys.
Lett. \textbf{22}, 145 (1993).
\bibitem{Giulini} D. Giulini, ``What is (not) wrong with scalar gravity?,'' arXiv:gr-qc/0611100.
\bibitem{dBS97} B. V. de Bakker and J. Smit, ``Gravitational binding in 4D dynamical triangulation,''
arXiv:hep-lat/9604023.
\bibitem{DLSU21} M. Dai, J. Laiho, M. Schiffer and J. Unmuth-Yockey, ``Newtonian binding from lattice quantum gravity,''
arXiv:2102.04492.
\end{thebibliography}

\end{document}
